\section{This project: what is proved here, and how the benchmark uses it}
\label{sec:repo}

\subsection{The heat-flux results}

The notes in \texttt{shared/docs/} and the manuscript
\texttt{plasma\_heat\_flux\_paper/manuscript.tex} study the function
$\chi(k,\omega)$ of Section~\ref{sec:between} for linear perturbations of a
Maxwellian, mostly at zero frequency. The manuscript is the current
statement; \texttt{summary\_of\_results.tex} predates several of its
sections. What follows is a reading guide, with the manuscript's
conditions. All of it is linear, for a homogeneous background, one Fourier
direction, and an unbounded or periodic domain. None of the proofs has been
checked by a second reader, and the Landau and electron results rest on
published coercivity estimates that are cited and not reproved.

\paragraph{Conventions.} The manuscript uses the energy flux, which in one
dimension is half the third central moment used in
Sections~\ref{sec:collisionless} and~\ref{sec:between}. Its
one-dimensional collisionless constant is $\sqrt{2/\pi}$, and its
conductivities are half those quoted above. Relative deficits are
unaffected (Appendix~\ref{app:conv}).

\paragraph{The response depends on the drive.} A response ratio is tied to
its input. Two drives are treated: a potential force (the Hammett--Perkins
setting) and an isotropic heating source. They share the conductivity and
differ at order $k^2$.

\begin{center}
\renewcommand{\arraystretch}{1.3}
\begin{tabularx}{\textwidth}{@{}P{3.7cm}YP{3.7cm}@{}}
\toprule
Setting & Established & Not established \\
\midrule
One velocity dimension, potential drive, any conserving symmetric
non-negative operator
  & dual variational principles;
    $\kF/(1+Kk^2)\le\chi(k)\le\kF$ at every $k$; $K$ is the exact leading
    deficit (under stated solvability hypotheses)
  & nonlinear or arbitrary-initial-data closure \\
BGK and Dougherty
  & closed forms \eqref{eq:theta-bgk}, \eqref{eq:theta-dougherty}; both
    limits; sharp leading deficits; exact frequency response for
    constant-$\nu$ BGK
  & equivalence to Coulomb collisions \\
Three dimensions, potential drive
  & a stress cross term appears; two-sided bounds at small $k$; leading
    coefficient; positivity when $k^2\kF\eta<3$
  & bounds at all $k$: they are false in general \\
Counterexamples
  & relaxation operators with a slow stress have $\chi>\kF$; with a slow
    heat flux as well, $\chi<0$
  & sharp thresholds \\
Like-particle Landau
  & weighted solvability; $|k|\chi\to3\sqrt{2/\pi}$ with rate
    $O(|k|^{-2/3}\log|k|)$; positivity at small and at large $k$
  & positivity at intermediate $k$; $\chi\le\kF$ at all $k$ \\
Three dimensions, thermal drive
  & $\kappa/(1+K_{\mathrm{th}}k^2)\le\chi_{\mathrm{th}}\le\kappa$ at every
    $k$, without momentum conservation; $\chi_{\mathrm{th}}$ is a Stieltjes
    function of $k^2$
  & the spectral measure beyond two moments \\
Electrons with ions of charge $Z$
  & the thermal theorems with the zero-current conductivity;
    $|k|\chi_{\mathrm{th}}\to\tfrac95\sqrt{2/\pi}$ for every $Z$; the
    normalized deficit grows linearly in $Z$ with slope $1.4305$
  & certified constants; a rate for the collisionless limit \\
Frequency dependence
  & the exact closure for data in the retained moments is a convolution in
    time with a positive-real kernel; bounds on the error of the
    quasi-static closure
  & the continuum statement for Landau \\
\bottomrule
\end{tabularx}
\end{center}

\paragraph{Numbers.} From polynomial Galerkin calculations: for
like-particle Landau collisions $\kF/\tau\simeq3.9502$ and
$K/\tau^2\simeq22.35$, a leading deficit about $7.45$ times that of the
Dougherty operator at matched conductivity. For electrons the normalized
deficit $K_{\mathrm{th}}/\kappa^2$ grows from about $4.3$ at $Z=1$ to $96$
at $Z=64$. In the embedded tests the exact quasi-static closure has
transient errors of up to $34\%$ in a one-dimensional model coupled to
sound waves, and stays within $2\%$ in the quasi-neutral decay of an
electron temperature perturbation at $Z=1$ and $|k|\lambda_e\le0.3$, where
Spitzer--H\"arm conduction is off by $42\%$.

\paragraph{What they are for.} These are analytic test targets for
classical and learned closures, with known answers at both ends and proved
intervals in between. They are not a general constitutive law.

\subsection{The benchmark tracks, in the five items}

\begin{center}
\renewcommand{\arraystretch}{1.3}
\small
\begin{tabularx}{\textwidth}{@{}P{2.3cm}P{2.6cm}P{2.5cm}P{2.5cm}Y@{}}
\toprule
Track & Hierarchy and limit & Retained & Discarded & Defect measured by \\
\midrule
A1 particles to Boltzmann
  & BBGKY; low density
  & one-particle distribution
  & cumulants, collision history
  & factorization error, recollision counts, Boltzmann residual \\
A2 waves to wave kinetics
  & correlation hierarchy; weak nonlinearity
  & spectrum
  & phases, spectral cumulants
  & spectral covariance error, kinetic residual \\
A3 Boltzmann to fluid
  & moment hierarchy; $\Kn\to0$
  & $n$, $\bu$, $T$
  & non-Maxwellian velocity structure
  & constitutive flux error, Hilbert remainder \\
B1 dusty plasma
  & three stacked closures (Section~\ref{sec:strong}); no small parameter
    assumed
  & grains; then $f_1$; then moments
  & plasma and wakes; correlations; velocity detail
  & force decomposition, factorization defect, flux residuals \\
B2 supernova
  & Reynolds decomposition of 3D onto 1D; none
  & radial profiles
  & nonradial motion, history
  & $C^\star_{\mathrm{SN}}$ by term \\
\bottomrule
\end{tabularx}
\end{center}
Track A is where the closure is a theorem and the asymptotic-limit error
can be separated from everything else. Track B is where the assumptions
fail and an augmented closure is needed. The heat-flux results sit between
them: a kinetic-to-moment closure with exact answers and a regime
parameter ($k\lambda$, and for electrons $Z$) that sweeps from the trusted
limit to its breakdown.

\subsection{The closure models compared}

The list in \texttt{mathematical\_closure\_program.md} maps onto the
sections of this tutorial as follows.
\begin{center}
\renewcommand{\arraystretch}{1.2}
\begin{tabularx}{\textwidth}{@{}YP{4.6cm}@{}}
\toprule
Model in the program & Section here \\
\midrule
Analytic and identifiable effective interaction laws & \ref{sec:strong} \\
Contextual or finite-memory graph particle models
  & \ref{sec:strong}, \ref{sec:memory} \\
Classical Markov kinetic or fluid closure
  & \ref{sec:asym}, \ref{sec:plasma}, \ref{sec:collisionless} \\
Explicit finite-cumulant or finite-moment closure
  & \ref{sec:hier}, \ref{sec:moment} \\
Finite-memory closure with kernels or auxiliary variables
  & \ref{sec:memory}, \ref{sec:between} \\
Structure-preserving learned correction & \ref{sec:learned} \\
Black-box surrogate as a control & \ref{sec:learned} \\
\bottomrule
\end{tabularx}
\end{center}
